\documentclass[12pt]{article}
\usepackage[utf8]{inputenc}
\usepackage[T1]{fontenc}
\usepackage{hyperref}
\title{Goreng Language Resources for K-12 Curriculum Development}
\author{Compiled by Claude Sonnet 5.5 (claude-sonnet-5-5) for Andr\'as Kornai}
\date{}
\begin{document}
\maketitle

\section{Introduction}
Goreng (Koreng; ISO 639-3 xgg) is treated as a variety of Noongar in south-western Australia, centred near Albany. Gerhardt Laves recorded Koreng in 1930, and this material was lost for many years and has only recently been recovered. I found no Goreng-specific teaching materials. I could not confirm any link to Ravensthorpe from the results.

\section{Primary Educational Resources}
No Goreng primer, textbook or school programme was found. Noongar resources exist more widely, but none from the searches were Goreng-specific.

\section{Linguistic and Reference Materials}
\begin{itemize}
\item Wikipedia, ``Noongar language'', which lists Koreng (Goreng) as xgg under the Noongar macrolanguage: \url{https://en.wikipedia.org/wiki/Noongar_language}
\item Bindon and Chadwick, \emph{A Nyoongar wordlist from the south-west of Western Australia}, which compiles several earlier word lists: \url{https://link.stirling.wa.gov.au/resource/k8mq-i6z1xE} and \url{https://link.belmont.wa.gov.au/resource/MiktFl81xzc}. These are library catalogue records.
\item AIATSIS Noongar (W41) entry, listing Thieberger (1993) among its sources: \url{https://collection.aiatsis.gov.au/node/1510236}
\end{itemize}

\section{Religious and Community Resources}
No Goreng-specific scripture or community organisation was found.

\section{Digital and Audio Resources}
No Goreng-specific audio, video or software was found.

\section{Curriculum Development Context}
Goreng material is sparse and mostly sits inside Noongar sources. A curriculum would need to work with Noongar community language programmes and draw on the Laves and Douglas material, checking regional variation carefully. The wordlists above are the starting points.
\end{document}